\section{Videre arbeid}
\label{sec:videre-arbeid}

Videre arbeid bør ta utgangspunkt i både testavgrensningene og
sluttevalueringen fra HDO.Dette er nødvendig for å gå fra en fungerende proof-of-concept til en løsning som kan vurderes for reell drift. Samtidig var den samlede
vurderingen 4 av 5, og én respondent presiserte at løsningen krever noe
videreutvikling før videre bruk.

Det viktigste neste steget er produksjonsherding av sikkerheten.
Issuer- og audience-validering bør aktiveres, \texttt{RequireHttpsMetadata}
bør være på, og den midlertidige Keycloak-konfigurasjonen bør erstattes
av HDOs faktiske autentiseringsoppsett. Tilgangsmodellen bør også vurderes
på nytt dersom flere klienter eller roller skal bruke samme API.

Testoppsettet bør utvides med sikkerhetsskanning i CI, lasttesting med
målbare resultater og kontraktstesting mot AntMedia, NHN og eventuelle nye
videomotorer. Dette er nødvendig for å gå fra en fungerende
proof-of-concept til en løsning som kan vurderes for reell drift.

Observability-oppsettet bør tilpasses HDOs driftsmiljø. \texttt{/metrics},
Prometheus og Grafana bør sikres med nettverksbegrensning, reverse proxy
eller annen tilgangskontroll, og dashboard, alerting og loggkorrelasjon
bør settes opp etter HDOs egne rutiner.

Til slutt bør nye providere behandles som egne integrasjonsprosjekter.
Valideringen viste at arkitekturen kan utvides uten endringer i
kjernelogikken, men hver produksjonsklar provider trenger egne tester,
dokumentasjon, health checks, metrikker og tydelig håndtering av
operasjoner som ikke støttes.